\documentclass[11pt]{article}
\usepackage[T1]{fontenc}
\usepackage{lmodern}
\usepackage[margin=1in]{geometry}
\usepackage{microtype}
\usepackage{amsmath,amssymb,amsthm,mathtools}
\usepackage[colorlinks=true,linkcolor=blue,citecolor=blue,urlcolor=blue]{hyperref}
\hypersetup{pdftitle={The Matrix Multiplication Bound omega <= 9/4 over Associative Rings},pdfauthor={},pdfsubject={AI-generated manuscript on an integral extension of the OpenAI matrix multiplication proof}}
\ifdefined\pdfinfoomitdate\pdfinfoomitdate=1\fi
\ifdefined\pdftrailerid\pdftrailerid{}\fi
\newtheorem{theorem}{Theorem}[section]
\newtheorem{lemma}[theorem]{Lemma}
\newtheorem{proposition}[theorem]{Proposition}
\theoremstyle{remark}
\newtheorem{remark}[theorem]{Remark}
\DeclareMathOperator{\rk}{R}
\newcommand{\Z}{\mathbb Z}
\newcommand{\Q}{\mathbb Q}
\newcommand{\C}{\mathbb C}
\newcommand{\F}{\mathbb F}
\newcommand{\eps}{\varepsilon}
\newcommand{\res}{\preceq}
\newcommand{\tp}[1]{^{\otimes #1}}
\numberwithin{equation}{section}
\title{The Matrix Multiplication Bound\\$\omega\leq9/4$ over Associative Rings}
\author{}
\date{}
\begin{document}
\maketitle

\begin{abstract}
OpenAI's auxiliary-separation argument gives the matrix multiplication
exponent bound $9/4$ over the complex numbers \cite{OAI,OAILean}; our earlier
extension establishes the same bound over every field \cite{Fields}.
We describe an integral version of that argument. Its stronger
intermediate conclusion is a bound on exact tensor rank over $\Z$, which
supplies the same exponent upper bound over every associative unital ring.
The zero ring is included using an explicit zero-cost program and the
real-infimum convention in the formal definition.
The changes are a primitive tensor semiring for spectral detection,
coefficient extraction and finite free descent, and coprime patching of
unnormalized Fourier and interpolation constructions. The numerical
exponent and the underlying profile-growth argument are those of OpenAI.
\end{abstract}

\noindent\textbf{Verification status.} The full upper-bound theorem,
including the zero ring, has compiled and passed its axiom audit,
fresh-kernel replay, and frozen-specification Comparator with all bundled checkers.
These checks validate formal declarations, not this prose document.

\paragraph{AI disclosure.}
This manuscript is AI-generated. The proof development and Lean
formalization were produced with Codex under human direction. Attribution
to the underlying OpenAI proof and the previous field extension is given
below; no claim of independent numerical improvement or human peer review
is made.

\section{Statement and contribution}

For an associative unital ring $R$, define $\omega(R)$ using division-free
arithmetic programs: addition, subtraction, and multiplication each cost
one operation; input and constant gates cost zero. The order of factors
in a multiplication gate is significant. An exponent $\tau$ is admissible
when
\[
 \forall\eps>0\ \exists C>0\ \forall n\geq1\ \exists P_n:\quad
 P_n\text{ computes matrix multiplication over }R,\qquad
 \operatorname{cost}(P_n)\leq Cn^{\tau+\eps}.
\]
Correctness means correctness for every pair of input matrices of the
specified size. The exponent is the real-valued \texttt{sInf} of admissible
exponents, with the convention for the zero ring explained below.
This is the arithmetic model used in the OpenAI Lean development, with
its scalar type generalized from a field to a ring.

Write $M_n$ for the square matrix multiplication tensor and $\rk_\Z(T)$
for exact tensor rank with integer coefficient vectors. Put
\[
 \nu_\Z=\inf_{n\geq2}\log_n\rk_\Z(M_n).
\]
A tensor identity here is equality at every coefficient, not merely
agreement of the corresponding functions over finite fields.

\begin{theorem}\label{thm:main}
The exact integer coefficient rank exponent satisfies $\nu_\Z\leq9/4$.
Consequently, for every associative unital ring $R$ and every $\eps>0$,
there are division-free arithmetic programs computing square matrix
multiplication with cost $O_{R,\eps}(n^{9/4+\eps})$.
For every such $R$, the formally defined exponent satisfies
\[
 \omega(R)\leq\frac94.
\]
\end{theorem}

\begin{remark}[The zero ring and the real infimum]\label{rem:zero-ring}
In the zero ring, one free constant-zero register supplies every output,
giving a correct program of cost zero. Every real exponent is therefore
admissible. The admissible set is $\mathbb R$, which is not bounded below;
Mathlib's real-valued \texttt{sInf} convention gives $\omega(R)=0$.
The extended-real infimum would instead be $-\infty$.
For nontrivial rings, admissible exponents are bounded below by $2$, so the
usual infimum interpretation applies. The integral argument handles that
case; the zero-ring case uses the explicit program and the stated convention.
\end{remark}

OpenAI supplies the complex-field bound, auxiliary separation,
determinant and sector constructions, and the real profile-growth
argument \cite{OAI,OAILean}. The previous field extension changes the
Fourier period in positive characteristic and descends from an algebraic
closure through a fixed finite extension \cite{Fields}. We do not infer
Theorem~\ref{thm:main} from the existence of solutions over all fields.
Instead we construct and patch integer coefficient identities. This also
handles rings with zero divisors, nilpotents, and noncommutative
multiplication. The coefficient algebras used in the proof remain
commutative.

\section{The integral spectral setting}

Write $U\res T$ if three integer linear maps restrict $T$ to $U$.
Pass to tensors modulo mutual restriction, with direct sum as addition
and tensor product as multiplication. Let $\mathcal S$ be the subsemiring
consisting of zero and the classes satisfying $1\res T$.
We call these classes \emph{unit-dominating}. A tensor with a coefficient
$1$ or $-1$ is unit-dominating.

This restriction on the spectral setting matters. The nonzero scalar
tensor with coefficient $2$ does not restrict over $\Z$ to the scalar
tensor $1$: such a restriction would give $2abc=1$ for integers $a,b,c$.
Thus nonzero tensors cannot simply be treated as positive spectral
objects in the field proof.

For nonzero $T\in\mathcal S$ one has
\[
 1\res T\res\langle\rk_\Z(T)\rangle,
\]
where $\langle m\rangle$ denotes the diagonal tensor of $m$ copies of
$1$. These comparisons give precisely the cofinality and positivity used
by the detecting-character argument of \cite[Appendix~A]{OAI}.
The order of diagonal tensors is the ordinary order of their sizes,
as follows from flattening over $\Q$. Flattening gives $\nu_\Z\geq2$; the naive multiplication scheme gives
$\nu_\Z\leq3$.

The detecting-character proof therefore runs on $\mathcal S$: for
$d\geq2$ and an integer $0\leq k<d^{\nu_\Z}$, there is a normalized,
additive, multiplicative, restriction-monotone character $\varphi$ with
$\varphi(M_d)\geq k$. Its compactness, rational separation, and fixed-point
steps are inherited from \cite{OAILean}; every division by a character
value is justified by $\varphi(T)\geq1$ for nonzero $T\in\mathcal S$.
The catalytic rank argument extracts $1$ only from such an object.

For convenience, these characters extend to all tensor classes:
\begin{equation}\label{eq:extend}
 \chi(T)=\varphi(T\oplus1)-1.
\end{equation}
Since $T\oplus1$ is unit-dominating, this is defined everywhere.
Additivity and multiplicativity follow by adjoining copies of $1$ and
expanding $(T\oplus1)\otimes(U\oplus1)$; monotonicity and nonnegativity
follow from restriction. The extension agrees with $\varphi$ on
$\mathcal S$. In particular it preserves the detecting value on $M_d$.
This avoids asserting positivity for every nonzero integer tensor.

\section{Transport of exact coefficient identities}

The following two elementary constructions replace field-valued
interpolation in the treatment of polynomial degenerations.

\begin{lemma}[Fixed finite free descent]\label{lem:descent}
Let $E$ be a commutative $\Z$-algebra free of finite rank $D$, with a
$\Z$-linear map $\pi:E\to\Z$ satisfying $\pi(1)=1$.
If integer tensors $U,T$ satisfy $U\res_E T$, then
\[
 U\res_\Z \bigoplus_{D^2}T.
\]
The same overhead $D^2$ applies after taking any tensor power in this
one fixed algebra $E$.
\end{lemma}

\begin{proof}
Choose a basis $e_1,\ldots,e_D$ of $E$. Expand the first two local
maps as $A=\sum_i A_i e_i$ and $B=\sum_j B_j e_j$, with integer
coefficient maps $A_i,B_j$. For each pair $(i,j)$ apply
$c\mapsto\pi(e_i e_jc)$ to the third local map. Summing the resulting
$D^2$ integer restrictions recovers $U$, because $\pi$ splits integer
coefficients. Taking powers before this expansion retains $D^2$;
descending first and then taking powers would instead incur $D^{2k}$.
\end{proof}

\begin{lemma}[Coefficient extraction]\label{lem:extract}
Suppose polynomial local maps $A(t),B(t),C(t)$ of degrees at most
$a,b,c$ transform $T$ into a polynomial tensor whose coefficient of
$t^h$ is $U$. Then $U$ is a restriction of $(a+1)(b+1)$ copies of $T$.
For a degeneration with leading coefficient $U$, its $k$th tensor power
has extraction cost at most $(ka+1)(kb+1)$.
\end{lemma}

\begin{proof}
The required coefficient is the sum of the restrictions with maps
$A_i,B_j,C_{h-i-j}$ for $0\leq i\leq a$ and $0\leq j\leq b$.
A negative coefficient index contributes zero. Powers multiply the
leading coefficient and scale each degree bound by $k$.
\end{proof}

Both lemmas hold over commutative base rings with the corresponding
finite free data. No evaluations or division are needed. Monic quotients
$E=\Z[t]/(f)$ supply the needed map $\pi$: reduce to degree below $\deg f$
and take the constant coefficient, assuming $\deg f>0$.

We also use a literal Bezout patch. If $d_1,\ldots,d_s$ generate the
unit ideal in $\Z$, their $k$th powers do too. Choose integers $u_j$ with
$\sum_j u_jd_j^k=1$. Concatenating rank decompositions for $d_j^kT\tp{k}$,
and multiplying one local map by $u_j$, gives a decomposition for
$T\tp{k}$. This is a coefficient identity; it does not divide by any
$d_j$ or treat scalar multiplication as a number of direct-sum copies.

\section{Integral convolution and Fourier separation}

\subsection{Interpolation with coprime norms}

For positive integers $a,b$, let $C(a,b)$ be the polynomial multiplication tensor with coefficients
$[i+j=\ell]$, and put $h=a+b-1$. The required statement is an asymptotic
rank bound, not a rank-$h$ decomposition over $\Z$ at the first power.

\begin{proposition}\label{prop:convolution}
For each $a,b$ there is a positive integer $K_{a,b}$ such that
\[
 \rk_\Z\bigl(C(a,b)\tp{k}\bigr)\leq K_{a,b}h^k
 \qquad(k\geq0).
\]
Every normalized nonnegative character monotone under restriction obeys
$\chi(C(a,b))\leq h$.
\end{proposition}

\begin{proof}
Fix a prime $p$. Choose a finite extension of $\F_p$ with at least $h$
elements and a power basis, then lift its monic minimal polynomial to a
monic $f\in\Z[t]$. The quotient $E=\Z[t]/(f)$ is finite free of some
positive rank $D$, and its basis specializes to the finite-field basis.
Choose lifts $v_1,\ldots,v_h$ of distinct residue elements. Write $V$ for
their Vandermonde matrix and $\Delta=\det V$.

Evaluation followed by $\operatorname{adj}(V)$ gives an $h$-term scheme
for $\Delta C(a,b)$ over $E$. Let $M_\Delta$ be the integer matrix of
multiplication by $\Delta$ in the chosen basis, and put $d=\det M_\Delta$.
If $w$ is the element with coordinate vector
$\operatorname{adj}(M_\Delta)[1]$, then $\Delta w=d\,1$.
Multiplying the output map by $w$ therefore gives an $h$-term scheme
for $dC(a,b)$ over $E$, without passing to a localization.

Modulo $p$, the nodes remain distinct. Thus $\Delta$ is nonzero in the
residue field, multiplication by it is invertible, and $p\nmid d$.
Taking powers in this fixed $E$ and applying
Lemma~\ref{lem:descent} gives
\[
 \rk_\Z\bigl(d^k C(a,b)\tp{k}\bigr)\leq D^2h^k.
\]
Perform this construction for each prime. No prime divides all resulting
clearing integers, so the ideal they generate is $\Z$. Membership of $1$
in that ideal uses only a finite subfamily $d_1,\ldots,d_s$.
Fix this subfamily and its algebras before varying $k$. The Bezout patch
gives the assertion with $K_{a,b}=\sum_jD_j^2$. Taking $k$th roots of
the character inequality removes this fixed constant.
\end{proof}

\subsection{Two unnormalized Fourier periods}

For $M\geq1$ matched sectors, take $L=5M$ and use each period
$q\in\{L,L+1\}$. Work in the finite free cyclotomic algebra
$E_q=\Z[z]/(\Phi_q(z))$, with $\zeta$ the image of $z$. It is a unit
and satisfies the unnormalized identity
\[
 \sum_{r=0}^{q-1}\zeta^{re}=
 \begin{cases}q,&q\mid e,\\0,&q\nmid e.\end{cases}
\]
Apply the substitutions of \cite[Proposition~3.1]{OAI}, omitting their
factor $1/q$. The relevant phases have absolute value at most
$3(M-1)<q$, so the same matched terms survive. The relative monomial
weights select the same leading tensor, multiplied by $q$. To make all
local exponents nonnegative, add $M^2$ to each second- and third-leg
weight. The selected coefficient moves to degree $2M^2$; its tensor is
unchanged.

Let $A$ be the shared-first-leg source and $B$ its separated target,
including the auxiliary dot-product factor. The construction gives a
polynomial degeneration from $q$ copies of $A$ to the coefficient
multiple $qB$ over $E_q$. Taking powers first, then extracting and
descending, gives
\[
 q^kB\tp{k}\res_\Z
 \bigoplus_{D_q^2 P_q(k)q^k} A\tp{k},
\]
where $D_q=\deg\Phi_q$ and $P_q$ is a fixed quadratic polynomial bound.
The $q^k$ on the left multiplies coefficients; on the right it counts
copies. Since $L^k$ and $(L+1)^k$ are coprime, patch the two restrictions
before taking character values. Then
\[
 \chi(B)^k\leq
 \bigl(D_L^2P_L(k)L^k+D_{L+1}^2P_{L+1}(k)(L+1)^k\bigr)\chi(A)^k.
\]
Taking roots yields $\chi(B)\leq(L+1)\chi(A)\leq6M\chi(A)$.
The constant six replaces the original constant five and disappears in
the entropy limit. For each fixed outer type length, the sector count
$M$, both algebras, and all degree bounds are fixed before the inner
power $k$ tends to infinity. No estimate uniform in growing $M$ is needed.

\section{The retained profile argument and ring algorithms}

The determinant filtration uses an integral adapted basis. On its affine
chart, for pairs $(A(X),B(X))$ of degree at most $e$, take quotient
vectors $(X^i,0)$ for $0\leq i\leq e$ and $(0,X^e)$, and kernel vectors
$(-X^{j+1},X^j)$ for $0\leq j<e$. The inverse reads quotient coordinates
from $A+XB$ and kernel coordinates from $B_0,\ldots,B_{e-1}$.
These are mutually inverse integer maps; no field dimension argument is
needed. The sector construction similarly uses coordinate maps and
integer monomial weights.

Thus the determinant, shifted-tripling, and entropy inequalities of
\cite{OAI} retain their form, while Proposition~\ref{prop:convolution}
supplies the needed convolution upper bound. For the symmetrized profile
$P$ and its parameter $t>0$, the retained laws include symmetry,
row concavity, shifted tripling
\[
 P(a,3b+a-1)\geq3P(a,b),
 \qquad P(a,b)\leq(a+b-1)^{1/t}.
\]
OpenAI's real diagonal-growth argument forces $t\leq3/4$, hence
$\chi(M_n)\leq n^{9/4}$. Detecting characters from Section~2 then give
$\nu_\Z\leq9/4$. This step invokes their proof rather than supplying a
new profile inequality or a better numerical exponent.

Fix a block size $u\geq2$ and an integer rank-$r$ scheme with
$r<u^{9/4+\delta}$, where $0<\delta<\eps$. Every output of the scheme is
\[
 \sum_{\ell=1}^r w_{c\ell}
 \left(\sum_a u_{\ell a}x_a\right)
 \left(\sum_b v_{\ell b}y_b\right),
 \qquad u_{\ell a},v_{\ell b},w_{c\ell}\in\Z.
\]
Integer constants are central in every associative unital ring.
Expansion always keeps $x_a$ before $y_b$, so the coefficient identity
works without commuting input variables. It also works for matrix
blocks, enabling the original recursive program construction.
Its recurrence is $S(uv)\leq rS(v)+6ru^2v^2$ with $u,r$ fixed.
Using $r\geq u^2$, padding, and positive exponent slack gives the
operation bound in Theorem~\ref{thm:main}. The direct operation-count
statement covers every ring. Nontriviality remains a hypothesis of the
bounded-below and lower-bound statements; the zero-ring exponent is handled
as in Remark~\ref{rem:zero-ring}.

\section{Formal scope and provenance}

The exported theorem in namespace
\texttt{OAI.MatrixMultiplication} is
\begin{verbatim}
theorem omega_le_nine_quarters_all_rings
    (R : Type u) [Ring R] :
    Arithmetic.omega R <= (9 : Real) / 4
\end{verbatim}
Lean's \texttt{Ring} class here means associative unital ring, not
necessarily commutative or nontrivial. The bounded-below admissibility and
$\omega(R)\geq2$ theorems retain \texttt{[Nontrivial R]}; neither holds
for the zero ring. Auxiliary tensor rank and coefficient-algebra
constructions are over commutative rings, with the main intermediate
specialized to $\Z$.

The formal conclusion is existence of a family of arithmetic programs.
It does not assert an $O(n^{9/4})$ endpoint bound, an efficient uniform
circuit generator, or bit complexity for encodings of ring elements.
In particular an arbitrary ring is not assumed to have a finite-tape
representation. The integer schemes are stronger coefficient statements
than functional correctness modulo a chosen characteristic.

The new repository \cite{Rings} starts from the minimal import closure
of the prior all-fields entry point at revision
\texttt{08481ef22bca7dc9ffba091083b7c1e81e537220}. Its import manifest
records the inherited source hashes and OpenAI upstream revision.
The current toolchain is Lean 4.35.0-rc3 with Mathlib revision
\texttt{c55e6e786f49471c72fbddbec5415808896aec1e}
\cite{Lean,Mathlib}.

Within the matrix-multiplication subtree, the relevant files under
\nolinkurl{AuxiliarySeparation/} are:
\begin{itemize}
\item \nolinkurl{Tensor/Primitive*.lean} and
\nolinkurl{Character/Existence.lean}: the spectral domain and detection;
\item \nolinkurl{Integral/}: coefficient extraction,
finite free descent, norm clearing, coprime patching, Fourier separation,
and the arithmetic bridge;
\item \nolinkurl{Determinant/} and
\nolinkurl{Sector/}: the integral versions of the inherited constructions.
\end{itemize}
The checked transport declarations include
\begin{quote}\small\ttfamily
convolution\_integer\_power\_rank\\
exists\_prime\_clearing\_convolution
\end{quote}
and the polynomial degeneration \texttt{power}/\texttt{extract} methods.

The six audited claims use only \texttt{propext}, \texttt{Classical.choice},
and \texttt{Quot.sound}. Their complete dependencies passed fresh-kernel
replay: the upper bound for all rings, boundedness and nonemptiness of the
admissible set, the lower bound, direct costs, and integer coefficient
identities. Boundedness and the lower bound retain nontriviality hypotheses.
Comparator checks these against a frozen model with exactly seven
field-to-ring substitutions and no definition holes. Its Lean paranoid,
lean4lean, nanoda, con-leche, con-ron, and default Lean checkers all accepted
the strengthened statement.
Verification uses trusted local sources, pinned dependencies, and incremental
build caches. Comparator's build sandbox is disabled. The repository records
commands, outcomes, prior failure-control tests, and separate AI source reviews.

\paragraph{Attribution.}
The numerical $9/4$ bound and its main mathematical architecture are due
to OpenAI \cite{OAI,OAILean}. The prior field extension is credited in
\cite{Fields}. This manuscript concerns the integral spectral setting,
explicit multiplier elimination, fixed-algebra transport, and the passage
to arbitrary associative rings. Standard descent, adjugate, and Bezout
principles are used without a historical-priority claim. The manuscript
and the extension's Lean proof development are AI-generated; neither is
presented as independent human peer review.

\begin{thebibliography}{9}
\raggedright
\bibitem{OAI}
OpenAI, \emph{An Upper Bound of $9/4$ for the Matrix Multiplication
Exponent}, October 2, 2026.
\href{https://github.com/openai/math/blob/adc7f1241b42e322a6451854ab7e4b4c146bf78a/preprints/Matrix-Multiplication-Nine-Fourths-October-2-2026/paper.pdf}{Preprint},
revision \texttt{adc7f1241b42e322a6451854ab7e4b4c146bf78a}.

\bibitem{OAILean}
OpenAI, \emph{Lean sources for matrix multiplication}, 2026.
\href{https://github.com/openai/math/tree/adc7f1241b42e322a6451854ab7e4b4c146bf78a/lean/OAI/LinearAlgebra/MatrixMultiplication}{MatrixMultiplication source subtree},
same pinned revision.

\bibitem{Fields}
S.~Navot, \emph{Matrix multiplication exponent at most $9/4$ over
arbitrary fields: a Lean formalization}, version 1.0.0, 2026.
\href{https://doi.org/10.5281/zenodo.23219128}{doi:10.5281/zenodo.23219128}.
The present source baseline is
\href{https://github.com/selanavot/matrix-multiplication-all-fields/tree/08481ef22bca7dc9ffba091083b7c1e81e537220}{revision 08481ef2}
of that repository; see its provenance records for the archived release.

\bibitem{Rings}
\emph{Matrix multiplication exponent at most $9/4$ over associative
rings: Lean development}, source repository.
\url{https://github.com/selanavot/omega_le_nine_quarters_all_rings}.
Source, paper, import provenance, and verification records are versioned
together; the manuscript is associated with its repository commit.

\bibitem{Lean}
L.~de Moura and S.~Ullrich, \emph{The Lean 4 theorem prover and
programming language}, Automated Deduction---CADE~28,
LNCS \textbf{12699} (2021), 625--635.
\href{https://doi.org/10.1007/978-3-030-79876-5_37}{doi:10.1007/978-3-030-79876-5\_37}.

\bibitem{Mathlib}
The mathlib Community, \emph{The Lean mathematical library},
CPP~2020, 367--381.
\href{https://doi.org/10.1145/3372885.3373824}{doi:10.1145/3372885.3373824}.
\end{thebibliography}
\end{document}
